\chapter{LaTeX Hints}
\label{sec:latexhints}

% Required for proper example rendering in the compiled PDF
\newcount\LTGbeginlineexample
\newcount\LTGendlineexample
\newenvironment{ltgexample}%
{\LTGbeginlineexample=\numexpr\inputlineno+1\relax}%
{\LTGendlineexample=\numexpr\inputlineno-1\relax%
  \tcbinputlisting{%
    listing only,
    listing file=\currfilepath,
    colback=green!5!white,
    colframe=green!25,
    coltitle=black!90,
    coltext=black!90,
    left=8mm,
    title=Corresponding \LaTeX{} code of \texttt{\currfilepath},
    listing options={
        frame=none,
        language={[LaTeX]TeX},
        escapeinside={},
        firstline=\the\LTGbeginlineexample,
        lastline=\the\LTGendlineexample,
        firstnumber=\the\LTGbeginlineexample,
        basewidth=.5em,
        aboveskip=0mm,
        belowskip=0mm,
        numbers=left,
        xleftmargin=0mm,
        numberstyle=\tiny,
        numbersep=8pt%
      }
  }
}%

This chapter contains hints on writing LaTeX.
It focuses on minimal examples, which can be directly adapted to the content

\section{Handling of paragraphs}

\begin{ltgexample}
One sentence per line.
This rule is important for the usage of version control systems.
A new line is generated with a blank line.
As you would do in Word:
New paragraphs are generated by pressing enter.
In LaTeX, this does not lead to a new paragraph as LaTeX joins subsequent lines.
In case you want a new paragraph, just press enter twice!
This leads to an empty line.
In word, there is the functionality to press shift and enter.
This leads to a hard line break.
The text starts at the beginning of a new line.
In LaTeX, you can do that by using two backslashes (\textbackslash\textbackslash).
\\
This is rarely used.

Please do \textit{not} use two backslashes for new paragraphs.
For instance, this sentence belongs to the same paragraph, whereas the last one started a new one.
A long motivation for that is provided at \url{http://loopspace.mathforge.org/HowDidIDoThat/TeX/VCS/#section.3}.
\end{ltgexample}

\section{Notes separated from the text}

The package mindflow enables writing down notes and annotations in a way so that they are separated from the main text.

\begin{ltgexample}
\begin{mindflow}
This is a small note.
\end{mindflow}
\end{ltgexample}

\section{Handling TODOs}

\begin{ltgexample}
\textmarker{Markierter Text.}
\end{ltgexample}

Bei \verb1\textmarker1 wird nur die Textfarbe geändert, da dies auch bei einigen Worten gut funktioniert.

\begin{ltgexample}
\textcomment{Markierter Text.}{Kommentar dazu.}
\end{ltgexample}

\begin{ltgexample}
\hl{In Gelb hervorgehoben.}
Provided indirectly by pdfcomment.sty (soulpos).
\end{ltgexample}

\begin{ltgexample}
\modified{Manuelle Markierung für Text, der seit der letzten Version geändert wurde.}
\end{ltgexample}

\begin{ltgexample}
Das ist ein Text.
\change{FL1: Text angepasst}{Geänderter Text}.
\end{ltgexample}

\begin{ltgexample}
Hier nur ein Kommentar\sidecomment{Kommentar}.
\end{ltgexample}

\begin{ltgexample}
\todo{Hier muss noch kräftig Text produziert werden}
\end{ltgexample}

\section{Hyphenation}

\LaTeX{} automatically hyphenates words.
When using \href{https://ctan.org/pkg/microtype}{microtype}, there should be fewer hyphenations than in other settings.
It might be necessary to tweak the hyphenations nevertheless.
Here are some hints:

\begin{ltgexample}
In case you write \enquote{application-specific}, then the word will only be hyphenated at the dash.
You can also write \verb1applica\allowbreak{}tion-specific1 (result: applica\allowbreak{}tion-specific), but this is much more effort.

You can now write words containing hyphens which are hyphenated at other places in the word.
For instance, \verb1application"=specific1 gets application"=specific.
This is enabled by an additional configuration of the babel package.
\end{ltgexample}

\section{Typesetting Units}

\begin{ltgexample}
Numbers can be written plain text (such as 100), by using the \href{https://ctan.org/pkg/siunitx}{siunitx} package as follows:
\SI{100}{\km\per\hour},
or by using plain \LaTeX{} (and math mode):
$100 \frac{\mathit{km}}{h}$.
\end{ltgexample}

\begin{ltgexample}
\SI{5}{\percent} of \SI{10}{kg}
\end{ltgexample}

\begin{ltgexample}
Numbers are automatically grouped: \num{123456}.
\end{ltgexample}

\section{Surrounding Text by Quotes}

\begin{ltgexample}
Please use the \enquote{enquote command} to quote something.
Quoting with "`quote"' or ``quote'' also works.

\end{ltgexample}

\section{Cleveref examples}
\label{sec:ex:cref}

Cleveref demonstration: Cref at beginning of sentence, cref in all other cases.

\begin{figure}
  \centering
  \includegraphics[width=.75\linewidth]{example-image-a}
  \caption{Example figure for cref demo}
  \label{fig:ex:cref}
\end{figure}

\begin{table}
  \centering
  \begin{tabular}{ll}
    \toprule
    Heading1 & Heading2 \\
    \midrule
    One      & Two      \\
    Thee     & Four     \\
    \bottomrule
  \end{tabular}
  \caption{Example table for cref demo}
  \label{tab:ex:cref}
\end{table}

\begin{ltgexample}
\Cref{fig:ex:cref} shows a simple fact, although \cref{fig:ex:cref} could also show something else.

\Cref{tab:ex:cref} shows a simple fact, although \cref{tab:ex:cref} could also show something else.

\Cref{sec:ex:cref} shows a simple fact, although \cref{sec:ex:cref} could also show something else.
\end{ltgexample}

\section{Figures}

\begin{ltgexample}
\Cref{fig:label} shows something interesting.

\begin{figure}
  \centering
  \includegraphics[width=.8\linewidth]{example-image-golden}
  \caption[Simple Figure]{
    Simple Figure.
    Based on \citet{mwe}.
  }
  \label{fig:label}
\end{figure}
\end{ltgexample}

\section{Sub Figures}

An example of two sub figures is shown in \cref{fig:two_sub_figures}.

\begin{ltgexample}
\begin{figure}[!b]
  \centering
  \subfloat[Case I]{\includegraphics[width=.4\linewidth]{example-image-a}%
    \label{fig:first_case}}
  \hfil
  \subfloat[Case II]{\includegraphics[width=.4\linewidth]{example-image-b}%
    \label{fig:second_case}}
  \caption{Example figure with two sub figures.}
  \label{fig:two_sub_figures}
\end{figure}
\end{ltgexample}

\section{Tables}

\begin{ltgexample}
\begin{table}
  \caption{Simple Table}
  \label{tab:simple}
  \centering
  \begin{tabular}{ll}
    \toprule
    Heading1 & Heading2 \\
    \midrule
    One      & Two      \\
    Thee     & Four     \\
    \bottomrule
  \end{tabular}
\end{table}
\end{ltgexample}

\begin{ltgexample}
% Source: https://tex.stackexchange.com/a/468994/9075
\begin{table}
  \caption{Table with diagonal line}
  \label{tab:diag}
  \begin{center}
    \begin{tabular}{|l|c|c|}
      \hline
      \diagbox[width=10em]{Diag \\Column Head I}{Diag Column\\Head II} & Second & Third \\
      \hline
       & foo & bar              \\
      \hline
    \end{tabular}
  \end{center}
\end{table}
\end{ltgexample}


\section{Source Code}

\begin{ltgexample}
\Cref{lst:XML} shows source code written in XML.
\Cref{line:comment} contains a comment.

\begin{lstlisting}[
  language=XML,
  caption={Example XML Listing},
  label={lst:XML}]
<listing name="example">
  <!-- comment --> (* \label{line:comment} *)
  <content>not interesting</content>
</listing>
\end{lstlisting}
\end{ltgexample}

One can also add \verb+float+ as parameter to have the listing floating.
\Cref{lst:flXML} shows the floating listing.

\begin{ltgexample}
\begin{lstlisting}[
  % one can adjust spacing here if required
  % aboveskip=2.5\baselineskip,
  % belowskip=-.8\baselineskip,
  float,
  language=XML,
  caption={Example XML listing -- placed as floating figure},
  label={lst:flXML}]
<listing name="example">
  Floating
</listing>
\end{lstlisting}
\end{ltgexample}

One can also typeset JSON as shown in \cref{lst:json}.

\begin{ltgexample}
\begin{lstlisting}[
  float,
  language=json,
  caption={Example JSON listing -- placed as floating figure},
  label={lst:json}]
{
  key: "value"
}
\end{lstlisting}
\end{ltgexample}

Java is also possible as shown in \cref{lst:java}.

\begin{ltgexample}
\begin{lstlisting}[
  caption={Example Java listing},
  label=lst:java,
  language=Java,
  float]
public class Hello {
    public static void main (String[] args) {
        System.out.println("Hello World!");
    }
}
\end{lstlisting}
\end{ltgexample}

\section{Itemization}

One can list items as follows:

\begin{ltgexample}
\begin{itemize}
  \item Item One
  \item Item Two
\end{itemize}
\end{ltgexample}

With the package paralist, one can create itemizations with lesser spacing:

\begin{ltgexample}
\begin{compactitem}
  \item Item One
  \item Item Two
\end{compactitem}
\end{ltgexample}

One can enumerate items as follows:

\begin{ltgexample}
\begin{enumerate}
  \item Item One
  \item Item Two
\end{enumerate}
\end{ltgexample}

With the package paralist, one can create enumerations with lesser spacing:

\begin{ltgexample}
\begin{compactenum}
  \item Item One
  \item Item Two
\end{compactenum}
\end{ltgexample}

With paralist, one can even have all items typeset after each other and have them clean in the TeX document:

\begin{ltgexample}
\begin{inparaenum}
  \item All these items...
  \item ...appear in one line
  \item This is enabled by the paralist package.
\end{inparaenum}
\end{ltgexample}

\section{Abbreviations}

With \verb+\gls{...}+ you can enter abbreviations, the first time you call it, the long form is used.
When reusing \verb+\gls{..}+ the short form is automatically displayed.
The abbreviation is also automatically inserted in the abbreviation list.
With \verb+\glspl{...}+ the plural form is used.
If you want the short form to appear directly at the first use, you can use \verb+\glsunset{..}+ to mark an abbreviation as already used.
The opposite is achieved with \verb+\glsreset{..}+.

Abbreviations are defined in \verb+\content\ausarbeitung.tex+ by means of \verb+\newacronym{...}{...}{...}+.

More information at: \url{https://ctan.org/pkg/bib2gls}.

\begin{ltgexample}
At the first pass the \gls{fr} was 5.
At the second pass was \gls{fr} 3.
The plural form can be seen here: \glspl{er}.
To demonstrate what the list of abbreviations looks like for longer description texts, \glspl{rdbms} must be mentioned here.

\gls{dante} is a local \TeX\ user group.
The German-speaking local \TeX\ user group is \gls{dante}.
A \gls{gp} is a medical doctor.
I went to my surgery to see the \gls{gp}.
\end{ltgexample}

\todo{Include difference between acronym and abbreviation - \url{https://tex.stackexchange.com/a/154566/9075}}

\section{Other Features}

\begin{ltgexample}
The words \enquote{workflow} and \enquote{dwarflike} can be copied from the PDF and pasted to a text file.
\end{ltgexample}

\begin{ltgexample}
The symbol for powerset is now correct: $\powerset$ and not a Weierstrass p ($\wp$).

$\powerset({1,2,3})$
\end{ltgexample}

\begin{ltgexample}
Brackets work as designed:
<test>
One can also input backticks in verbatim text: \verb|`test`|.
\end{ltgexample}


\section{Varioref examples}
\label{sec:ex:vref}

Varioref demonstration: Vref at beginning of sentence, vref in all other cases.

\begin{ltgexample}
\Vref{fig:ex:cref} shows a simple fact, although \vref{fig:ex:cref} could also show something else.

\Vref{tab:ex:cref} shows a simple fact, although \vref{tab:ex:cref} could also show something else.

\Vref{sec:ex:cref} shows a simple fact, although \vref{sec:ex:cref} could also show something else.
\end{ltgexample}

\section{Citations}

When referencing something from the bibliography file, it will automatically appear in the references section.
If a reference is not cited, it is not appearing there.

\begin{ltgexample}
Standard case: Citing indirectly citing something~\cite{mwe}.
In case one wants to name the author: \textcite{mwe} shows a minimal \LaTeX{} example.
\end{ltgexample}

Note that \texttt{\textbackslash textcite\{mwe\}} prints both the author and the reference to the bibliography entry.

Remember that you have to call \texttt{biber main-english} to generate the bibliography data for \texttt{lualatex}.
You will need to run \texttt{lualatex} twice to ensure that the page numbers are updated correctly.


In the bibliography, use \texttt{\textbackslash textsuperscript} for \enquote{st}, \enquote{nd}, \ldots:
E.g., \enquote{The 2\textsuperscript{nd} conference on examples}.
When you use \href{https://www.jabref.org}{JabRef}, you can use the clean up command to achieve that.
See \url{https://help.jabref.org/en/CleanupEntries} for an overview of the cleanup functionality.

\section{Miscellaneous Examles}
\label{ssec:example}

Referencetest: \Cref{ssec:example}, \cref{fig:Abbildung} und \cref{alg:example}.

\begin{ltgexample}
Checkmark: \dingcheck.
Crossmark: \dingcross.
\end{ltgexample}

\begin{figure}
  \missingfigure{}
  \caption{Abbildung}
  \label{fig:Abbildung}
\end{figure}

\begin{landscape}
  \begin{figure}
    \missingfigure{}
    \caption{Gedrehte Abbildung}
    \label{fig:AbbildungGedreht}
  \end{figure}
\end{landscape}

\subsection{Algorithmen}

\begin{algorithm}
  \caption{$algo$}
  \label{alg:example}
  \begin{algorithmic}[1]
    \State $a \gets 0$
    \State State 2\label{alg1:state2}
  \end{algorithmic}
\end{algorithm}

\begin{algorithm}
  \caption{Algorithmus 2}
  \label{alg:example2}
  \begin{algorithmic}[1]
    \State $a \gets 0$
    \State State 2\label{alg2:state2}
  \end{algorithmic}
\end{algorithm}

\Cref{alg:example} hat bereits einen Algorithmus gezeigt.
Test der Zeilenreferenzierung: Zeile~\ref{alg1:state2} (\cref{alg:example}) und Zeile~\ref{alg2:state2} (\cref{alg:example2}).

\subsection{Definitionen}
\begin{definition}[Title]
  \label{def:def1}
  Definition Text
\end{definition}

\Cref{def:def1} zeigt \ldots

\subsection{Aufzählungen}

\begin{enumerate}[label=\alph*)]
  \item a
  \item b
  \item c
  \item d
\end{enumerate}

Equivalent to paralist's inparaenum:
\begin{enumerate*}[label=\alph*)]
  \item a
  \item b
  \item c
  \item d
\end{enumerate*}

\begin{description}
  \item[first] Erstens
  \item[second] Zweitens
  \item[third] Drittens
\end{description}

\begin{description}
  \item[\texttt{first}] Erstens
  \item[\texttt{second}] Zweitens
  \item[\texttt{third}] Drittens
\end{description}

%works only if package enumitem is loaded
\begin{description}[font=\ttfamily]
  \item[first] Erstens
  \item[second] Zweitens
  \item[third] Drittens
\end{description}

\begin{description}[style=unboxed]
  \item[first label with a long description text breaking over one line. Enabled by enumitem package] Erstens
  \item[second] Zweitens
  \item[third] Drittens
\end{description}

\begin{Description}
  \item[first label with a long description text breaking over one line. Defined in template.tex] Erstens
  \item[second] Zweitens
  \item[third] Drittens
\end{Description}

\begin{itemize}
  \item Erstens
  \item Zweitens
  \item Drittens
\end{itemize}

Optionaler Parameter ändert den Marker, der vorangestellt ist.
Siehe \url{http://www.weinelt.de/latex/item.html}.
\begin{itemize}
  \item[A] Erstens
  \item[B] Zweitens
  \item[C] Drittens
\end{itemize}

Falsche Benutzung des optionalen Parameters wie folgt:
\begin{itemize}
  \item[first] Erstens
  \item[second] Zweitens
  \item[third] Drittens
\end{itemize}
 ist zu beachten, dass es sich bei Einbindung von \texttt{enumitem} anders verhält als bei \texttt{paralist}.

\subsection{fquote}

\begin{fquote}[T.\ Informatiker]
  Bis nächsten Freitag ist das Programm fertig.
\end{fquote}

\begin{gfquote}{T.\ Informatiker}
  Bis nächsten Freitag ist das Programm fertig.
\end{gfquote}


%%% ===============================================================================
